% Agglomerative hierarchical clustering (Ward). Subsection of "Clustering".

\subsection{Agglomerative hierarchical clustering}  % bottom-up merging, Ward linkage
\label{sec:clust_agglo}                  % \secref{sec:clust_agglo}

% \paragraph{Principle}                    % bottom-up construction of a hierarchy
Agglomerative clustering builds a hierarchy of nested groups from the bottom up
\cite{norouzi2024volpam}. At the start, every tree is a cluster of its own. At each step, the two
closest clusters are merged, and the procedure repeats until a single cluster contains all the
trees. The sequence of merges forms a binary tree, the \emph{dendrogram}, in which the height of
each junction shows how far apart the two merged clusters were. Cutting the dendrogram with a
horizontal line at some height gives a partition into $K$ clusters, so one run of the algorithm
yields the clustering for every $K$ and shows how small groups nest inside larger ones. This is
what makes hierarchical methods attractive for discovering intrinsic groupings whose number is
not known in advance \cite{norouzi2024volpam}.

\paragraph{Linkage}                      % how "closest" is defined: Ward
What ``closest'' means is set by the linkage. \citet{norouzi2024volpam} merge the two
clusters whose centroids are closest. We use Ward's linkage \cite{ward1963}, which merges the
pair of clusters whose union increases the total within-cluster sum of squares the least. For
two clusters $A$ and $B$ with centroids $\boldsymbol{\mu}_A$ and $\boldsymbol{\mu}_B$, this
increase is
\begin{equation}
  \Delta(A, B) = \frac{|A|\,|B|}{|A| + |B|}\,
                 \big\|\boldsymbol{\mu}_A - \boldsymbol{\mu}_B\big\|_2^2 .
  \label{eq:ward}                        % cost of merging A and B
\end{equation}
Like the centroid linkage, Ward's criterion compares cluster centres, but the factor in front of
the distance penalizes merging two large clusters, so it favours compact clusters of similar
size. Since it greedily minimizes the same within-cluster sum of squares as K-means
(Eq.~\eqref{eq:kmeans}), the two methods can be compared directly.

\paragraph{Choosing the number of clusters}  % elbow on the average intra-cluster distance
To choose where to cut the dendrogram, we follow the elbow criterion of
\citet{norouzi2024volpam}. For every number of clusters $k = 1, \dots, 19$, we cut the hierarchy
into $k$ clusters and compute the average intra-cluster distance, the mean distance between two
trees of the same cluster, averaged over the clusters:
\begin{equation}
  D_{\mathrm{avg}}(k) = \frac{1}{k} \sum_{l=1}^{k}
     \frac{1}{N_l^2} \sum_{u=1}^{N_l} \sum_{v=1}^{N_l}
     \big\|\mathbf{x}_{ul} - \mathbf{x}_{vl}\big\|_2 ,
  \label{eq:davg}                        % mean pairwise distance inside each cluster
\end{equation}
where $N_l$ is the number of trees in cluster $l$ and $\mathbf{x}_{ul}$ is the code of its
$u$-th tree. Splitting a cluster always makes its parts more compact, so $D_{\mathrm{avg}}$ tends
to decrease with $k$. The useful information is where it drops the most: the number of
clusters is chosen as the one with the largest relative reduction \cite{norouzi2024volpam},
\begin{equation}
  k^{*} = \arg\max_{k}\;
          \frac{D_{\mathrm{avg}}(k-1) - D_{\mathrm{avg}}(k)}{D_{\mathrm{avg}}(k-1)} .
  \label{eq:kstar}                       % largest relative drop of D_avg
\end{equation}
A large relative drop means that going from $k-1$ to $k$ clusters separates groups that are
genuinely different; after that point, further splits mostly cut coherent groups into pieces.
The elbow criterion identified $k^{\*}=4$ as the optimal number of clusters for both the registered and unregistered airway trees. Thus, in both representations, the hierarchy was cut at four clusters for the subsequent analysis.